\usepackage{textcomp}
\usepackage{lipsum}
\usepackage[ngerman,english]{babel}
\usepackage{blindtext}
\usepackage{booktabs}
\usepackage{microtype}
\usepackage{amsmath, amssymb, amsthm}
\usepackage{comment}
\usepackage{ragged2e} % Good for forcing alignements
\usepackage{array}
\usepackage{graphicx,color}
\usepackage{pdfpages}
\makeatletter
\define@key{Gin}{artifact}[]{}
\makeatother
\usepackage{multirow}
\usepackage[
    a5paper,
    twoside,
    inner=27mm,
    outer=18mm,
    top=17mm,
    bottom=17mm
]{geometry}
\usepackage{epsfig}
\usepackage{setspace}
\usepackage{latexsym}
\usepackage{mathrsfs}
\usepackage{tikz}
\usepackage[titletoc]{appendix}
\usepackage{alltt}
\graphicspath{ {images/} }
\usepackage{longtable}
\usepackage{acronym}
\usepackage{pgfplots, pgfplotstable}
\usepackage{csquotes}
\usepackage{caption}
\usepackage{subcaption}
\usepackage[hidelinks]{hyperref}
\usepackage{algorithmicx}
\usepackage{cleveref } % for nice ref stuff
%\usepackage{tocloft}
%\setcounter{tocdepth}{1}
%\renewcommand{\cftdot}{}
\usepackage{titletoc,tocloft}
%\setlength{\cftsecindent}{2cm}
\setlength{\cftsubsecindent}{1.2cm}
\setlength{\cftsubsubsecindent}{1cm}
%\dottedcontents{section}[1.5em]{}{1.3em}{.6em}
%\usepgfplotslibrary{groupplots} % Not necessary

\usepackage{times}
% Line spacing - Distance between paragraph lines
\renewcommand{\baselinestretch}{1.28}
\emergencystretch=3em
\doublehyphendemerits=100000
\providecommand{\chapterfont}[1]{}
\providecommand{\sectionfont}[1]{}
\usepackage{titlesec}

\setcounter{secnumdepth}{3}
\chapterfont{\centering }
\titleformat{\chapter}[display]
{\normalfont%
\fontsize{11pt}{13pt}\selectfont% %change this size to your needs for the first line
\bfseries}{\chaptertitlename\ \thechapter}{14pt}{%
    \fontsize{12pt}{14pt}\selectfont %change this size to your needs for the second line
}
%{\bf\centering }
%{\chaptertitlename\ \thechapter }{14pt}{\large}

\usepackage{textcase}
%\allsectionsfont{\MakeTextUppercase}

\sectionfont{\normalfont}
\titleformat{\section}
{\normalfont\bfseries\fontsize{10pt}{12pt}\selectfont}{\thesection}{0.8em}{}
\titleformat{\subsection}
{\normalfont\bfseries\fontsize{9pt}{11pt}\selectfont}{\thesubsection}{0.8em}{}
\titleformat{\subsubsection}
{\normalfont\fontsize{8.5pt}{10.5pt}\selectfont}{\thesubsubsection}{0.8em}{}
\titlespacing*{\chapter}{0pt}{0pt}{2ex plus .3ex}
\titlespacing*{\section}{0pt}{2.4ex plus .3ex}{1ex plus .2ex}
\titlespacing*{\subsection}{0pt}{2ex plus .2ex}{0.8ex plus .1ex}
\titlespacing*{\subsubsection}{0pt}{1.6ex plus .2ex}{0.6ex plus .1ex}
%\titlespacing*{\section} {0pt}{0pt}{2.3ex plus .2ex}
%\titlespacing*{\subsection}{0pt}{0pt}{1.5ex plus .2ex}
%%%% CAPTION SETUP AND FONT SIZE %%%%%%%%%%%%%%%%%%%%%%%%%%%%%%%%%%%%%%%%%%%%%%%
\captionsetup[table]{labelsep=space, labelfont=bf, font=footnotesize, skip=4pt}
% Available values are scriptsize, footnotesize, small, normalsize, large,
% and Larger.
%
\captionsetup[figure]{labelsep=space, labelfont=bf, font=footnotesize, skip=4pt}

\newcommand{\printfitgraphics}[2][0.62\textheight]{%
    \includegraphics[width=\linewidth,height=#1,keepaspectratio]{#2}%
}

\setcounter{topnumber}{3}
\setcounter{bottomnumber}{2}
\setcounter{totalnumber}{4}
\renewcommand{\topfraction}{0.86}
\renewcommand{\bottomfraction}{0.72}
\renewcommand{\textfraction}{0.12}
\renewcommand{\floatpagefraction}{0.68}

\setcounter{page}{1}
\newtheorem{definition}{Definition}
\newcommand{\brdefinition}{\begin{definition}}
\newcommand{\erdefinition}{\end{definition}}

\newtheorem{corollary}{Corollary}
\newcommand{\bcorollary}{\begin{corollary}}
\newcommand{\ecorollary}{\end{corollary}}
\newtheorem{example}{Example}
\newcommand{\bexample}{\begin{example}}
\newcommand{\eexample}{\end{example}}
\newtheorem{remark}{Remark}
\newcommand{\bremark}{\begin{remark}}
\newcommand{\eremark}{\end{remark}}
\newcommand{\bproof}{{\bf {Proof:}}}
\newcommand{\eproof}{}
\newcommand{\bsolution}{{\bf {Solution:}}}
\newcommand{\esolution}{}
\newtheorem{theorem}{Theorem}
\newcommand{\btheorem}{\begin{theorem}}
\newcommand{\etheorem}{\end{theorem}}
\newtheorem{lemma}{Lemma}
\newcommand{\blemma}{\begin{lemma}}
\newcommand{\elemma}{\end{lemma}}
\newcommand{\UD}{\ensuremath{\bigtriangleup}}
\newcommand{\IC}{\ensuremath{\mathcal{C}} }
%\renewcommand{\rm}{\normalshape}
\newcommand{\ol}{\overline}

%%%% Paragraph Settings %%%%%%%%%%%%%%%%%%%%%%%%%%%%%%%%%%%%%%%%%%%%%%%%%%%%%%%%
\raggedbottom
\clubpenalty=10000
\widowpenalty=10000
\displaywidowpenalty=10000

\def\cen{\centerline}
\def\pnq{\par\noindent\quad}
\def\pn{\par\noindent}
% Space between paragrahs
\usepackage[skip=6pt plus1pt minus1pt, indent=28pt]{parskip}

\newcommand{\non}{\nonumber}
\newcommand{\mC}{{\mathbb C}}
\newcommand{\mN}{{\mathbb N}}
\newcommand{\mU}{{\mathbb U}}
\newcommand{\cA}{{\mathcal A}}
\newcommand{\cR}{{\mathcal R}}
\newcommand{\cS}{{\mathcal S}}
\newcommand{\cT}{{\mathcal T}}
\newcommand{\cUCV}{{\mathcal UCV}}
\newcommand{\cST}{{\mathcal ST}}
\newcommand{\cK}{{\mathcal K}}
\newcommand{\ds}{\displaystyle}
\newcommand{\brdef}{\begin{defi}}
\newcommand{\erdef}{\end{defi}}
\newcommand{\bcor}{\begin{cor}}
\newcommand{\ecor}{\end{cor}}
\newcommand{\bth}{\begin{thm}}
\newcommand{\ble}{\begin{lem}}
\newcommand{\ele}{\end{lem}}
\newcommand{\bcha}{\end{cha}}\pagestyle{plain}
\renewcommand{\theequation}{\thechapter.\arabic{equation}}
\renewcommand{\thetheorem}{\thesection.\arabic{theorem}}
\renewcommand{\thecorollary}{\thesection.\arabic{corollary}}
\renewcommand{\thelemma}{\thesection.\arabic{lemma}}
\renewcommand{\thedefinition}{\thesection.\arabic{definition}}
\renewcommand{\theexample}{\thesection.\arabic{example}}
\renewcommand{\theremark}{\thesection.\arabic{remark}}
\renewcommand{\thechapter}{\arabic{chapter}}
\renewcommand{\thesection}{\thechapter.\arabic{section}}

\renewcommand{\chaptername}{CHAPTER}
\renewcommand{\figurename}{Fig.}

\def\cen{\centerline}
\def\pnq{\par\noindent\quad}
\def\pn{\par\noindent}
\def\sevenpoint{%
    \def\rm{\sevenrm}%
    \def\it{\sevenit}%
    \def\bf{\sevenbf}%
    \rm}
%\pretoler
\renewcommand{\theequation}{\thechapter.\arabic{equation}}
\theoremstyle{definition}
\renewcommand*{\proofname}{{\rm Proof}}
% Chapter title format
\usepackage{titlesec}
\chapterfont{\centering }
\titleformat{\chapter}[display]
{\bf\centering}
{\chaptertitlename\ \thechapter}{12pt}{\fontsize{12pt}{14pt}\selectfont}
\sectionfont{\normalfont}
%\renewcommand*\contentsname{\large \centerline{TABLE OF CONTENTS}}
\renewcommand\cftchappresnum{} % prefix "Chapter " to chapter number in ToC
\cftsetindents{chapter}{0em}{3em}      % set amount of indenting
\cftsetindents{section}{0em}{3em}
%For reference

%\usepackage{natbib}
%%\bibliographystyle{apalike}%\usepackage{apalike}
%
%\setlength{\bibhang}{0pt}
%
%%\setlength\bibindent{0em}
%\renewcommand{\bibname}{REFERENCES}
%\renewcommand{\bibname}{\protect\centerline{References}}

%\newcommand{\setbibname}[1]{\bibname[#1]{\centering #1}}


\renewcommand*\contentsname{Table of Contents}
\renewcommand*\listfigurename{List of Figures}
\renewcommand*\listtablename{List of Tables}
\newcommand{\frontmatterlisttitle}[1]{%
    \vspace*{\cftbeforelottitleskip}%
    \noindent\makebox[\textwidth][l]{{\cftlottitlefont #1}}\par%
    \vskip\cftafterlottitleskip}

%For Terms and Abbreviations 
%\usepackage[acronym,toc]{glossaries}
\usepackage{glossaries}
\makeglossary
%\makeglossaries
%\renewcommand*\glspostdescription{\cftdotfill{\cftsecdotsep}}
%\renewcommand{\glossarysection}[2][]{{\centering\bfseries\MakeTextUppercase{#2}\par}}


%\renewcommand*{\bibpagespunct}{\addcolon\space}

\definecolor{butter1}{rgb}{0.988,0.914,0.310}
\definecolor{chocolate1}{rgb}{0.914,0.725,0.431}
\definecolor{chameleon1}{rgb}{0.541,0.886,0.204}
\definecolor{skyblue1}{rgb}{0.447,0.624,0.812}
\definecolor{plum1}{rgb}{0.678,0.498,0.659}
\definecolor{scarletred1}{rgb}{0.937,0.161,0.161}


%For Table Column Width
%\usepackage{array} %% Already included above
\newcolumntype{L}[1]
{>{\raggedright\let\newline\\\arraybackslash\hspace{0pt}}m{#1}}
\newcolumntype{C}[1]
{>{\centering\let\newline\\\arraybackslash\hspace{0pt}}m{#1}}
\newcolumntype{R}[1]
{>{\raggedleft\let\newline\\\arraybackslash\hspace{0pt}}m{#1}}

\usepackage{etoolbox}

\makeatletter
\patchcmd{\ttlh@hang}{\parindent\z@}{\parindent\z@\leavevmode}{}{}
\patchcmd{\ttlh@hang}{\noindent}{}{}{}
\makeatother

% For fonts
\renewcommand{\familydefault}{\sfdefault}
%\usepackage{bookman}
%\usepackage{lmodern}
%\usepackage[T1]{fontenc}

%%%% Date and Time %%%%%%%%%%%%%%%%%%%%%%%%%%%%%%%%%%%%%%%%%%%%%%%%%%%%%%%%%%%%%
\usepackage{datetime}
\newdateformat{monthdayyeardate}{%
    \THEDAY~\monthname[\THEMONTH]~\THEYEAR}

%%%%% Bibliography %%%%%%%%%%%%%%%%%%%%%%%%%%%%%%%%%%%%%%%%%%%%%%%%%%%%%%%%%%%%%
%\usepackage{biblatex} %Imports biblatex package

%\addbibresource{backmatter/mypubs.bib} %Import the bibliography file


%%%% Environments %%%%%%%%%%%%%%%%%%%%%%%%%%%%%%%%%%%%%%%%%%%%%%%%%%%%%%%%%%%%%%
% For defining symbols under equationns
% https://tex.stackexchange.com/questions/95838
\usepackage{calc, tabularx}
\newenvironment{definesymbolsfrontmatter}[1][{~}]
{\noindent#1 \small\begin{tabular}[t]{>{$}l<{$} @{${}\quad{}$} l}}
{\end{tabular}\par}

\newenvironment{definesymbols}[1][{\textbf{\scriptsize~Where:}}]
{#1 \scriptsize\begin{tabular}[t]{>{$}l<{$} @{${}={}$} l}}
{\end{tabular}\\}

% For forcing a paraghraph to move to bottom of a page
% https://tex.stackexchange.com/questions/31186
\newenvironment{bottompar}{\par\vspace*{\fill}}{\clearpage}

%%%% FIX FOOTNOTE SPLIT ACROSS PAGES %%%%%%%%%%%%%%%%%%%%%%%%%%%%%%%%%%%%%%%%%%%
\interfootnotelinepenalty=10000

%%%%% SI Units
\usepackage{siunitx}
\sisetup{per-mode = symbol}
\DeclareSIUnit{\belmilliwatt}{Bm}
\DeclareSIUnit{\dBm}{\deci\belmilliwatt}
\DeclareSIUnit{\dB}{\deci\bel}
\DeclareSIUnit{\ppm}{ppm}
\DeclareSIUnit{\kbd}{kBd}


%%%% DIFF Equations
\usepackage{derivative}
\usepackage{physics}
\newcommand*\diff{\mathop{}\!\mathrm{d}}
\newcommand*\Diff[1]{\mathop{}\!\mathrm{d^#1}}

%\pdv{f}{x}, \quad \odv{Q}{t}=\odv{s}{t}, \quad \pdv{f}{x,y}, \quad
%\derivset{\odv}[switch-*=false] \odv{y}{x}, \quad \odv[n]{y}{x}, \quad
%\derivset{\odv}[misc-add-delims=fun] \odv*{\odv{y}{x}}{x}, \quad
%\derivset{\pdv}[sort-method={sign,symbol,abs}] \pdv[c+kn,-b+2a]{f}{x,y}

%%%% For graphs and flowcharts
\usepackage{pgf-pie}
\usetikzlibrary{shadows, shapes.geometric, arrows}
\pgfplotsset{compat=1.18} % use package is declared in the top of this file
\newlength{\TL}  % Tick Label


% Define styles for nodes
\tikzstyle{startstop} = [rectangle, rounded corners, minimum width=1cm, minimum height=1cm, text centered, draw=black, fill=red!30, font=\scriptsize]
\tikzstyle{io} = [
trapezium,
trapezium stretches=true, % A later addition
trapezium left angle=70,
trapezium right angle=110,
, minimum width=1cm, minimum height=1cm, text centered, draw=black, fill=red!30, font=\scriptsize]
\tikzstyle{process} = [rectangle, minimum width=1cm, minimum height=1cm, text centered, draw=black, fill=orange!30, font=\scriptsize]
\tikzstyle{decision} = [diamond, minimum width=1cm, minimum height=1cm, text centered, draw=black, fill=green!30, aspect=2, font=\scriptsize]
\tikzstyle{arrow} = [thick,->,>=stealth]

%%%%% For tables
%\usepackage[utf8]{inputenc}
%\usepackage{fourier}
\usepackage{makecell}
%\renewcommand\theadalign{bc}
%\renewcommand\theadfont{\bfseries}
%\renewcommand\theadgape{\Gape[4pt]}
%\renewcommand\cellgape{\Gape[4pt]}
%\newcommand{\specialcell}[2][c]{%
%\begin{tabular}[#1]{@{}c@{}}#2\end{tabular}}

% For centering columns in a table
% Taken from
% https://tex.stackexchange.com/questions/157389
\newcolumntype{P}[1]{>{\centering\arraybackslash}p{#1}}
\newcolumntype{M}[1]{>{\centering\arraybackslash}m{#1}}

% Afterpage package to allow joinging things
% https://tex.stackexchange.com/questions/65606/
\usepackage{afterpage}


% Allows for adjusting line spacing and etc.
% https://stackoverflow.com/questions/38144458
\usepackage{enumitem}
% used like this:
% \begin{itemize}[topsep=8pt,itemsep=4pt,partopsep=4pt, parsep=4pt]

% For itemize pro/con lists
% https://tex.stackexchange.com/questions/145198
\newcommand\itempro{\item[$+$]}
\newcommand\itemcon{\item[$-$]}


%% For table with color
%% !!! THIS CAUSED OPTION CLASH!!!
%\usepackage[table]{xcolor}


% For floatbarrier
\usepackage{placeins}


% FOR CODE LISTINGS

\usepackage{listings}

\lstdefinestyle{commonstyle}{
keywordstyle=[1]\bfseries,
commentstyle=\itshape,
stringstyle=\ttfamily,
basicstyle=\ttfamily\scriptsize,
numbers=left,
numberstyle=\tiny,
stepnumber=1,
numbersep=5pt,
tabsize=2,
showspaces=false,
showstringspaces=false,
breaklines=true,
breakatwhitespace=true,
escapechar=@
}

% Define VHDL language settings
\lstdefinelanguage{VHDL}{
morekeywords=[1]{architecture,begin,end,entity,is,port,in,out,of,use,library,IEEE,std_logic_vector,std_logic,signal,if,else,case,when,others,process,and,or,not},
sensitive=false,
morecomment=[l]--,
morecomment=[s]{/*}{*/},
morestring=[b]",
morestring=[b]',
keywordstyle=[1]\bfseries,
commentstyle=\itshape,
stringstyle=\ttfamily,
basicstyle=\ttfamily\scriptsize,
numbers=left,
numberstyle=\tiny,
stepnumber=1,
numbersep=5pt,
tabsize=2,
showspaces=false,
showstringspaces=false,
breaklines=true,
breakatwhitespace=true
}

% Set common settings
\lstset{
frame=single,% adds a frame around the code
frameround=tttt, % rounded corners
captionpos=t, % sets the caption-position to bottom
keepspaces=true, % keeps spaces in text, useful for keeping indentation of code
aboveskip=4pt,  % Space above each line
belowskip=1pt,  % Space below each line
lineskip=-1pt   % Further reduce spacing between lines
}



%%%%
%%%%%
% for combining two figures into one but keeping captions separat

% for empty pages
% https://tex.stackexchange.com/questions/36880/insert-a-blank-page-after-current-page


\usepackage{afterpage}

\newcommand\blankpage{%
    \null
    \thispagestyle{empty}%
    \newpage}

\makeatletter
\renewcommand{\cleardoublepage}{%
    \clearpage
    \if@twoside
        \ifodd\c@page
        \else
            \hbox{}%
            \thispagestyle{empty}%
            \newpage
        \fi
    \fi}
\makeatother
